% Generated draft: requires mathematical and source review.
\documentclass[11pt]{article}
\usepackage[margin=25mm]{geometry}
\usepackage{fontspec}
\setmainfont{Latin Modern Roman}
\usepackage{amsmath,amssymb,mathtools,unicode-math}
\setmathfont{Latin Modern Math}
\usepackage{enumitem,needspace,xurl,hyperref}
\setlength{\parindent}{0pt}
\setlength{\parskip}{5pt}
\setlength{\emergencystretch}{3em}
\newcommand{\N}{\mathbb{N}}
\newcommand{\Z}{\mathbb{Z}}
\newcommand{\Q}{\mathbb{Q}}
\newcommand{\R}{\mathbb{R}}
\newcommand{\C}{\mathbb{C}}
\newcommand{\F}{\mathbb{F}}
\newcommand{\A}{\mathbb{A}}
\newcommand{\PP}{\mathbb{P}}
\newcommand{\E}{\mathbb{E}}
\newcommand{\eps}{\varepsilon}
\newcommand{\dd}{\,\mathrm{d}}
\begin{document}
\section*{UnsolvedMath ID: 701199}
\textbf{Source problem code:} catalogue entry
\subsection{Definitions and mathematical statement}
Under the usual finite-rank interpretation, let \(n\) be an integer with \(n\geq 2\), and let \(F_n\) denote the free group of rank \(n\). For any homomorphism \(\alpha:F_n\to\operatorname{Aut}(F_n)\), define the semidirect product
\[
G_\alpha=F_n\rtimes_\alpha F_n,
\qquad
(x,u)(y,v)=\bigl(x\alpha(u)(y),uv\bigr).
\]
The first copy of \(F_n\) is the normal factor, and the second acts through \(\alpha\). The action is arbitrary, including the trivial action.

For a fixed group \(G_\alpha\), interpret its membership problem as the uniform subgroup membership problem: given a finite list of words \(h_1,\ldots,h_k\) and a word \(g\) in the generators of the two free factors, decide whether
\[
g\in\langle h_1,\ldots,h_k\rangle\leq G_\alpha.
\]
Here \(k\geq 0\), and the subgroup generated by the empty list is the trivial subgroup. Decidability means that an algorithm halts with the correct answer on every such input; the algorithm may depend on the fixed group \(G_\alpha\).

\textbf{Question (V. G. Bardakov, Kourovka Notebook 16.7, 2006).} Is this membership problem undecidable for every integer \(n\geq 2\) and every homomorphism \(\alpha:F_n\to\operatorname{Aut}(F_n)\)?

The supplied statement writes the product as \(F_n\leftthreetimes F_n\). The formulation above makes an action convention explicit. Finite rank and uniformity over finitely generated subgroups are conventional interpretations of the supplied notation and terminology, which the record does not expressly define. In particular, the question as formulated does not assert that every such group contains a single fixed finitely generated subgroup with undecidable membership.

The source states that an affirmative answer would imply an answer to Kourovka Notebook Problem 6.24; that problem's statement and the intended implication are not supplied. It also recalls decidability of the membership problem for \(F_n\), attributed to M. Hall (1949), and undecidability for \(F_n\times F_n\), attributed to K. A. Mikhailova (1958). These are retained as supplied background claims.
\subsection{Short English statement}
For every non-abelian free group of finite rank, and every action of one copy on another, must the resulting semidirect product have an undecidable subgroup membership problem? Here the input includes finite generators of a subgroup and an element whose membership is to be decided. The source labels the problem open; this label is supplied metadata and has not been independently verified.
\subsection{Sources}
\begin{itemize}
\item Local UnsolvedMath registry, supplied complete source record: permanent ID \(701199\), exact source code \texttt{catalogue entry}, title \emph{Kourovka Notebook Problem 16.7}. Upstream revision: \texttt{d371ee917ac9647fbb06cd6fc529deec6f3dda82}.
\item \emph{The Kourovka Notebook (Unsolved Problems in Group Theory)}, Problem 16.7, proposed by V. G. Bardakov in the 16th issue (2006). The supplied background also identifies the 21st issue edition (2026), arXiv identifier \texttt{1401.0300v48}. The registry describes the collection as Issues 1--21 (1965--2026), edited by E. I. Khukhro and V. D. Mazurov. These source details are supplied attribution, not evidence of independent consultation.
\item Supplied background attribution: M. Hall (1949), decidability of membership for \(F_n\); K. A. Mikhailova (1958), undecidability of membership for \(F_n\times F_n\). No publication titles or fuller bibliographic details are supplied for these attributions.
\item Supplied cross-reference: Kourovka Notebook Problem 6.24. Its statement is absent from the record.
\end{itemize}
\end{document}
